\chapter{Introduction}
\label{chapter: introduction}

\par In 1935 aerodynamicists from all over the world gathered in Rome to discuss the behavior of air when objects reach high speeds, notably, the change in air density by the phenomenon of compressibility \cite{TheBookOfAnderson}. At the usual low subsonic speeds of the time, it was assumed that the variation in density was negligible. However, air density starts to considerably change with higher velocities, leading to an additional variable that needs to be accounted for and coupling the momentum and energy equations. For speeds faster than sound, these density gradients become even more relevant with the formation of shock waves, which represent an almost discontinuous change in the fluid thermodynamic properties. The developments of this conference completely changed the scope of flow physics and gave rise to a new era of aerospace vehicles; it was the dawn of compressible aerodynamics.

\par The study of compressible flows through experimental methods is still fundamental to the modern aerospace sector \cite{TheBookOfAnderson}. Many applications such as transonic, supersonic, hypersonic vehicles, propulsion and the energy sector, still heavily rely on experimental data from practical setups to aid its design and validate Computational Fluid Dynamics (CFD) models \cite{Raffel2001}. These setups require a method to visualize and quantify the variation of density in the fluid. Such techniques should be quantitative, non-intrusive, capture a large field of view, have sufficient spatial resolution, sensitivity and if possible, be simple. To meet these requirements, optical flow measurement techniques have been developed based on the behavior of light through fluids, such as Shadowgraphy, Schlieren, Interferometry and most recently Background Oriented Schlieren (BOS).

\par These optical methods designed to visualize density gradients in compressible flows are based on light’s refraction. The phenomenon refers to the bending of light when it encounters differences in the medium it is traveling in. These differences are quantified by the medium refractive index $n$, which in fluids is dependent on the wavelength of light and the its thermochemical state. In gases, this dependency is given by the Gladstone-Dale equation, a simplification of the Lorentz-Lorentz equation, that makes explicit the proportionality between $n$ and the gas density \cite{Schmidt2025}. In nature, this phenomenon can be seen when temperature gradients lead to air density variations, creating visual distortions such as mirages and background shimmer. By controlling the medium variables in a laboratory setting, this distortion can be quantified to infer a specific property of the flow. This way, it is able to non-intrusively measure the fluid's density and potentially assess other thermodynamic properties \cite{Michalski2018}.

\par The first widespread optical density measurement techniques, namely shadowgraphy and schlieren, infer the density mainly qualitatively through analog equipment and require large optical elements. BOS improves on it by applying the same principles with modern tools. The technique applies image cross-correlation algorithms to quantify the distortion of a reference background and in doing so, it provides a powerful quantitative displacement field on a simpler setup \cite{Settles2001, Schmidt2025}. This improvement upon traditional schlieren made BOS a widely used technique since its introduction 25 years ago, nonetheless, limitations regarding sensitivity, spatial resolution and optical configuration remain.

\par Laser-Speckled Background Oriented Schlieren (SBOS) emerges as an alternative capable of overcoming several of the constraints posed by BOS. This new method employs a laser-speckled pattern as the reference, which is always in focus, allowing for dynamic positioning of the background image by adjusting the camera setup. This opens the possibility of unusual reference placements, such as in front of the test section or even inside of it, that might lead to higher sensitivity and spatial resolution compared to standard BOS configuration \cite{Meier&Roesgen2013}. Moreover, the high luminous intensity of the laser generated background allows for lower exposure times, making it appropriate for the imaging of fast phenomena. It additionally provides high luminous contrast, making it suitable for the visualization of light emitting low noise to signal phenomena, such as reactive flows \cite{Nakamura2021}. These advantages compared to standard BOS makes the study of the method interesting in an academic standpoint and potentially useful for a myriad of industrial applications.

\par However, the technique is still relatively unexplored compared to classical BOS systems and lacks a coherent theoretical framework for its implementation. Key aspects such as optimal setup design guidelines, uncertainty quantification, the implementation to complex high-speed wind tunnels and the implications of unusual background placements have yet to be thoroughly addressed in the literature. As will be established in the theoretical review of Chapter \ref{chapter: LiteratureReview}, these open questions represent a meaningful gap in the current state of the art, motivating the need for a structured investigation of the method and its application in new study cases.

Following the current research in compressible flows at TU Delft, this thesis will implement SBOS in the visualization of the flow through a linear aerospike nozzle in supersonic wind tunnel. Aerospike nozzles provide high adaptability to variations of ambient pressure during the rise in altitude of a launcher vehicle\cite{Wisse2005}. This characteristic makes them an interesting solution to vehicles designed to reach orbit in a single stage as it promises an increase in efficiency of up to 15\% compared to non-adaptive nozzles \cite{Hagemann1998}. SBOS's increased sensitivity and spatial resolution makes it a suitable method to visualize the intricacies present in the ramp-flow interaction of an aerospike. Current literature has yet to implement SBOS in a supersonic wind tunnel environment, making this experiment the first of its type. 

\par The research objective of this work is to study the compressible flow of an aerospike exhaust through the implementation of the SBOS method in a supersonic wind tunnel. Foundation and practical guidelines to future research will be created by designing, building, tuning and investigating uncertainties in a simple setup as proof of concept. The results will be used to verify the methodology. Later, the knowledge and experience gathered will be used to prepare the method on a more complex test stand, TU Delft's ST-15 high speed wind tunnel, to visualize the complex features present at an aerospike nozzle exhaust flow. To meet this objective, the following research questions are posed:
\begin{itemize}
    \item How can SBOS be systematically implemented to study compressible flows in a laboratory and wind tunnel setting?
    \begin{itemize}
        \item What should be the practical guidelines for the design of a SBOS setup, balancing the trade-offs between sensitivity and spatial resolution?
        \item  What are the dominant sources of uncertainty in SBOS measurements?
        \item How does SBOS compare in sensitivity, spatial resolution and accuracy to standard BOS?
        \item How does positioning the reference inside the test section affect the results' sensitivity and data treatment?
        \item What are the challenges and limitations of applying SBOS in a high-speed wind tunnel environment to visualize the flow of a linear aerospike? 
    \end{itemize}
\end{itemize}

\par The next chapters of this thesis are structured as follows: Chapter \ref{chapter: Theory} provides the theoretical background; Chapter \ref{chapter: LiteratureReview} provides a literature review; Chapter \ref{chapter: methodology} provides practical guidelines to design a BOS setup; Chapter \ref{chapter: Experimental Setups} explains the experiments conducted; Chapter \ref{chapter: Results} presents the results and data analysis; Chapter \ref{chapter: Conclusion} discusses main take aways of the experiments, summarizing the conclusions of the project.
